\chapter{Beyond markets}\label{ch:beyondmarkets}

Nothing in Chapters~\ref{ch:framework} and~\ref{ch:beyondlinear} is specific to prices. The framework needs a state that the agents watch and move, agents whose responses are filters of that state, and a safety question about how far the state can travel under bounded disturbances. This chapter maps three other systems into that template. One mapping is exact and its safety bound follows from Theorem~\ref{thm:sector}; the others are stated as targets, and all three empirical tests are pending.

\section{The template}\label{sec:bm_template}

A system fits the framework when (a) a scalar or vector state follows a stable linear plant driven by exogenous disturbances and by the agents' aggregate action; (b) each agent's action is a sector-bounded function of a filtered version of the state; and (c) the operator's safety question is a bound on the state's peak excursion under disturbances bounded by $\bar w$. Under (a)--(c), Theorem~\ref{thm:sector} and Proposition~\ref{prop:multiasset} give a safety bound that depends on the agents only through their total sensitivity, and Theorem~\ref{thm:dich} says what it costs to certify that sensitivity from experiments. Identification of the individual agents remains as hard as in Chapter~\ref{ch:p1}.

\section{Grid frequency with unknown device populations}\label{sec:bm_grid}

In a single-area power system, the frequency deviation $f$ obeys the swing equation $M\dot f=\Delta P-Df$, where $M$ is aggregate inertia, $D$ load damping and $\Delta P$ the power imbalance. Discretized with step $\Delta$, $f_{t+1}=(1-D\Delta/M)f_t+(\Delta/M)\Delta P_{t+1}$. Write the imbalance as an exogenous part $w$ (load steps, generator trips) plus the aggregate response of $K$ device classes (inverter-based resources, flexible loads) that react to filtered frequency, $u_t=\sum_iq_i\varphi_i(y_{i,t})$. This is the economy of Definition~\ref{def:cohort} with $a=1-D\Delta/M$ and $\kappa=\Delta/M$; devices that support frequency enter with negative weights, devices that amplify deviations with positive weights.

\begin{corollary}[Certified frequency excursion]\label{cor:grid}
In the discretized single-area model, if the device classes satisfy Assumption~\ref{ass:sector} with total sensitivity $\bar\Phi$ and $a+\kappa\bar\Phi<1$, the frequency deviation after any sequence of imbalances bounded by $\bar w$ never exceeds $\bar w\,\kappa/(1-a-\kappa\bar\Phi)$, whatever the devices' time constants, control laws within the sector bounds, or changes over time. \hfill\Proved
\end{corollary}

The corollary is Theorem~\ref{thm:sector} read in grid units. Its value is that an operator can bound worst-case frequency from an aggregate sensitivity, which can be certified from controlled perturbations or from responses to recorded disturbances, without device-level models of every aggregator's fleet. Three caveats matter. The bound counts every device class through $|q_i|$, so supporting devices loosen it rather than tighten it; when most devices support frequency the bound is conservative, and a sign-aware version is a \Target{}. Real systems have governor and automatic generation control dynamics, which add plant states; the multi-state version requires Proposition~\ref{prop:multiasset} or a sharper positive-systems argument (\Target{}). And grid-forming inverters respond to phase as well as frequency, which the scalar template does not represent.

\pendbox{grid test}{Using publicly available sub-second frequency recordings around documented imbalance events, estimate the aggregate frequency response, certify total sensitivity with the bootstrap procedure of Chapter~\ref{ch:p2}, and compare the certified excursion bound with realized nadirs. Data availability and the event list are to be confirmed before registration.}

\section{Vehicle platoons and string stability}\label{sec:bm_platoon}

In a platoon, each vehicle's controller responds to the spacing and speed of the vehicle ahead, and its own motion becomes the next vehicle's input. A platoon is string stable in the $\ell_\infty$ sense when disturbances do not grow as they travel down the string, which for linear controllers means that the $\ell_1$ norm of the error-propagation impulse response is at most one \citep{t:swaroop1996}. Commercial adaptive cruise controllers are proprietary, heterogeneous and updated over the air: they are unknown adaptive agents. Field experiments suggest that several are not string stable \citep{t:gunter2021}, and public data from car-following experiments with commercial systems exist \citep{t:makridis2021}.

The certification question, ``is this platoon string stable?'', depends on the controllers only through the error-propagation response, not through their internal structure. Under a declared bound on controller memory, the summed absolute response is certified from perturbation experiments at the rate of Theorem~\ref{thm:dich}(i) with $a_t\equiv1$ applied to the absolute taps, while recovering each controller's parameters is an identification problem of the kind studied in Chapter~\ref{ch:p1}.

\begin{proposition}[Certifying string stability without identifying controllers]\label{prop:platoon}
For a platoon of vehicles with unknown linear controllers whose error-propagation impulse responses $\gamma$ satisfy $|\gamma_t|\le C\bar r^{\,t}$, the $\ell_1$ string-stability margin $1-\|\gamma\|_1$ can be certified to accuracy $\epsilon$ from $O(\sigma^2\epsilon^{-2}\log(1/\epsilon))$ perturbation experiments, whatever the controllers' internal order, and certification is impossible without a bound of this kind. \hfill\Target
\end{proposition}

The upper bound follows the proof of Theorem~\ref{thm:dich}(i) for a non-linear but Lipschitz function of the taps; the impossibility part follows part (iii). The statement needs care with sign changes in $\gamma$, which make $\|\gamma\|_1$ a non-smooth functional; this is the technical content of the target.

\pendbox{platoon test}{Using recorded car-following experiments with commercial adaptive cruise control \citep{t:makridis2021}, certify the $\ell_1$ string-stability margin for each tested configuration and compare with the stability conclusions reported for the same systems. Access conditions and the list of usable experiments are to be confirmed before registration.}

\section{Markets of trading algorithms that learn}\label{sec:bm_ai}

Markets in which learning algorithms, including reinforcement-learning and language-model agents, trade with each other are the original setting of this thesis pushed forward. Two results apply directly. Certification of safety needs only a bound on total sensitivity (Theorem~\ref{thm:sector}); and the certificate lasts for its safety margin divided by the agents' learning speed (Corollary~\ref{thm:lifetime}). The second result gives a supervisory rule with no counterpart in identity-based regimes: recertify at a frequency set by how fast the agents update.

\pendbox{multi-agent simulation}{Simulate markets of reinforcement-learning trading agents in the price model of Chapter~\ref{ch:p2}, measure the realized lifetime of certificates as agents keep learning, and compare it with Corollary~\ref{thm:lifetime}. The specification (agents, learning rates, horizons and pass criteria) is to be registered before the simulation is run.}
